% !TEX program = xelatex
\documentclass[11pt,a4paper]{article}
\usepackage{fontspec}
\IfFontExistsTF{Georgia}{\setmainfont{Georgia}}{\setmainfont{TeX Gyre Termes}}
\usepackage[spanish,es-nodecimaldot]{babel}
\usepackage[margin=23mm]{geometry}
\usepackage{parskip,enumitem,xcolor,xurl,hyperref,fancyhdr,longtable,booktabs,needspace}
\newcommand{\UniversityName}{Universidad de Oviedo}
\newcommand{\SchoolName}{Escuela Politécnica de Ingeniería de Gijón}
\newcommand{\DegreeName}{Grado en Ciencia e Ingeniería de Datos}
\newcommand{\DocumentType}{Trabajo de Fin de Grado}
\newcommand{\ThesisTitle}{Análisis de ECGs con grandes modelos de lenguaje: aplicación a enfermedades raras}
\newcommand{\ThesisTitleEnglish}{ECG analysis with large language models: application to rare diseases}
\newcommand{\ThesisSubtitle}{Estudio comparativo del aprendizaje supervisado y el aprendizaje en contexto}
\newcommand{\ThesisRunningTitle}{Análisis de ECGs con grandes modelos de lenguaje}
\newcommand{\AuthorName}{Alejandro Moris Lara}
\newcommand{\SupervisorName}{Nahuel Costa Cortez}
\newcommand{\SubmissionPlace}{Gijón}
\newcommand{\SubmissionDate}{noviembre de 2026}
% Generated from the included canonical aggregate snapshot.
\newcommand{\SimCnnBA}{0,885}
\newcommand{\SimVlmBA}{0,904}
\newcommand{\TransferCnnBestBA}{0,551}
\newcommand{\TransferCnnLastBA}{0,498}
\newcommand{\TransferVlmBestBA}{0,514}
\newcommand{\RealVlmBA}{0,761}
\newcommand{\RealVlmSens}{0,741}
\newcommand{\RealVlmSpec}{0,781}
\newcommand{\RealVlmFone}{0,699}
\newcommand{\RealCnnBestBA}{0,546}
\newcommand{\RealCnnLastBA}{0,535}
\newcommand{\SimPositive}{20}
\newcommand{\SimNegative}{80}
\newcommand{\SimPatients}{100}
\newcommand{\HucaPositive}{116}
\newcommand{\HucaNegative}{201}
\newcommand{\HucaPatients}{317}

\hypersetup{hidelinks,pdftitle={Guion de defensa del TFG},pdfauthor={\AuthorName}}
\pagestyle{fancy}\fancyhf{}\fancyhead[L]{\small Guion de defensa del TFG}\fancyhead[R]{\small Alejandro Moris Lara}\fancyfoot[C]{\thepage}
\setlength{\headheight}{15pt}
\begin{document}
\begin{center}{\LARGE Guion para la defensa del TFG}\par\medskip
{\large\ThesisTitle}\par\medskip\AuthorName\par Noviembre de 2026\end{center}
\section*{Preparación y duración}
La defensa contiene 39 diapositivas principales para un ensayo de 18 minutos y dos páginas de referencias.

La apertura sitúa el problema cardiovascular y presenta la propuesta antes del índice. Después del estado del arte se explica cómo leer un latido. Tras las implementaciones se presentan los tres escenarios y el protocolo por paciente. Una diapositiva de transición abre cada bloque de resultados.

La referencia pública localizada para la EPI indica entre 15 y 20 minutos de forma orientativa y deja la duración al tribunal. El procedimiento del curso 2026--2027 está enlazado en SharePoint y requiere acceso institucional. Por tanto, estos tiempos son objetivos de ensayo, no una confirmación del límite de tu tribunal.

Ensaya en voz alta, siguiendo el gráfico que ve el tribunal. No leas las diapositivas ni recites los decimales. Al explicar una curva, identifica primero el eje horizontal, la métrica y las series, y después señala el cambio relevante. Las barras muestran variación entre semillas, no intervalos de confianza.

La presentación es una adaptación a 16:9 de los modelos del apartado 6.1 del Manual de Identidad Corporativa de UniOvi. No constituye una plantilla Beamer publicada o aprobada por la institución. Se conservan los emblemas originales, con Georgia para los títulos y el texto. Si Georgia no está instalada, el proyecto utiliza TeX Gyre Termes como alternativa.

\clearpage\section*{Ruta principal: 18 minutos}
\begin{longtable}{p{12mm}p{112mm}r}\toprule Diap. & Contenido & Acumulado\\\midrule\endhead
1 & Análisis de ECGs con grandes modelos de lenguaje & 0:15\\
2 & Enfermedades cardiovasculares y síndrome de Brugada & 0:40\\
3 & Electrocardiograma en formato imagen & 1:15\\
4 & Propuesta y objetivos del TFG & 1:45\\
5 & Estructura de la exposición & 2:00\\
6 & Estado del arte & 2:25\\
7 & Cómo leer un latido & 3:00\\
8 & Síndrome de Brugada y derivación V1 & 3:30\\
9 & Redes neuronales convolucionales & 3:50\\
10 & Aprendizaje en contexto & 4:15\\
11 & Adaptación de dominio & 4:35\\
12 & Conjuntos de datos del simulador y del HUCA & 5:00\\
13 & Simulador QRS/ST & 5:35\\
14 & Familias del simulador & 6:00\\
15 & Distribución de etiquetas en el simulador & 6:25\\
16 & Conjunto clínico del HUCA & 6:50\\
17 & Diseño experimental & 7:20\\
18 & Implementación de la CNN con ResNet18 & 7:55\\
19 & Inspección de activaciones con Grad-CAM & 8:25\\
20 & Implementación del ICL con Gemma 4 E4B & 8:50\\
21 & Escenarios de evaluación & 9:05\\
22 & Protocolo de evaluación por paciente: LOOCV & 9:40\\
23 & Simulación & 9:45\\
24 & Entrenamiento y evaluación con datos simulados & 10:40\\
25 & Matrices de confusión en el simulador & 11:10\\
26 & Cambio de dominio & 11:15\\
27 & CNN entrenada en el simulador y evaluada en el HUCA & 11:50\\
28 & Métodos de adaptación de dominio en la CNN & 12:30\\
29 & F1 de la clase positiva al cambiar de dominio & 13:00\\
30 & Cambio de dominio en ICL con selección balanceada & 13:35\\
31 & Cambio de dominio en ICL con selección estándar & 13:55\\
32 & Evaluación clínica & 14:00\\
33 & Entrenamiento y demostraciones con ejemplos del HUCA & 15:10\\
34 & Matrices de confusión en la evaluación clínica & 15:45\\
35 & Aspectos éticos y protección de datos & 16:00\\
36 & Discusión y limitaciones & 16:45\\
37 & Conclusiones & 17:25\\
38 & Continuación del trabajo & 17:50\\
39 & Cierre y preguntas & 18:00\\
\bottomrule\end{longtable}
\Needspace{8\baselineskip}\section*{01. Análisis de ECGs con grandes modelos de lenguaje}
\textbf{Tiempo orientativo: }15 segundos.

Buenos días. Este TFG estudia la clasificación de imágenes de ECG con pocos ejemplos etiquetados y utiliza el síndrome de Brugada como caso de estudio.

\textbf{Transición.} Empiezo por la disponibilidad de datos etiquetados.

\Needspace{8\baselineskip}\section*{02. Enfermedades cardiovasculares y síndrome de Brugada}
\textbf{Tiempo orientativo: }25 segundos.

Según la OMS, las enfermedades cardiovasculares causaron unos 19,8 millones de muertes en el mundo en 2022. Dentro de este grupo hay trastornos poco frecuentes, como el síndrome de Brugada. En Brugada cuesta reunir suficientes ECG etiquetados. El TFG explora cómo aprovecharlos para una futura herramienta de apoyo al triaje.

\textbf{Transición.} Además de cuántos ECG hay, importa en qué formato están disponibles.

\Needspace{8\baselineskip}\section*{03. Electrocardiograma en formato imagen}
\textbf{Tiempo orientativo: }35 segundos.

Esta fotografía muestra una forma en que puede conservarse o compartirse un ECG. En algunos registros solo está disponible el trazado impreso o su copia como imagen. Para analizarlo, una opción es reconstruir la serie temporal. Aquí estudiamos la clasificación directa a partir de imágenes del trazado. En el experimento, la entrada de ambos modelos es la imagen de un latido de V1.

\textbf{Transición.} A partir de esta representación, el TFG pregunta qué rendimiento podemos obtener con pocos ejemplos etiquetados.

\Needspace{8\baselineskip}\section*{04. Propuesta y objetivos del TFG}
\textbf{Tiempo orientativo: }30 segundos.

El TFG explora cómo apoyar el triaje de ECG sospechosos. En el análisis de ECG, este TFG estudia tres preguntas. ¿Cómo varía el rendimiento cuando se dispone de más ejemplos, ya sea para entrenar la red o para construir el contexto? ¿Cómo se compara una red convolucional supervisada con un modelo de visión y lenguaje que aprende en contexto? ¿En qué medida se conserva el rendimiento al pasar de ejemplos sintéticos a registros clínicos?

\textbf{Transición.} La exposición sigue el recorrido que aparece en el índice.

\Needspace{8\baselineskip}\section*{05. Estructura de la exposición}
\textbf{Tiempo orientativo: }15 segundos.

Primero situaré el trabajo y explicaré sus fundamentos. Después presentaré los datos, el simulador y los modelos. Antes de los resultados veremos los tres escenarios y cómo se evalúa cada paciente.

\textbf{Transición.} El estado del arte distingue tres formas de utilizar los ECG y sus etiquetas.

\Needspace{8\baselineskip}\section*{06. Estado del arte}
\textbf{Tiempo orientativo: }25 segundos.

Ribeiro y colaboradores estudiaron la clasificación de señales de ECG con registros de la red de telemedicina de Minas Gerais. Sangha y colaboradores entrenaron con registros brasileños convertidos en imágenes y realizaron una validación externa que incluyó PTB-XL. Camba y Costa estudiaron métodos con pocos ejemplos en ECG. Con ICL, los ejemplos etiquetados se incorporan directamente a la consulta.

\textbf{Transición.} Antes de ver el caso de Brugada, vamos a localizar las partes de un latido.

\Needspace{8\baselineskip}\section*{07. Cómo leer un latido}
\textbf{Tiempo orientativo: }35 segundos.

Leemos el trazado de izquierda a derecha. La P corresponde a la activación de las aurículas y el QRS, a la de los ventrículos. En el bloqueo de rama derecha, o RBBB, el ventrículo derecho se activa con retraso y puede aparecer una segunda deflexión positiva al final del QRS. En el simulador, RBBB etiqueta ese rasgo morfológico. Después del punto J miramos si el ST se eleva respecto a la línea de base y si la onda T, que corresponde a la repolarización ventricular, apunta hacia abajo.

\textbf{Transición.} Veamos ahora cómo aparecen estos rasgos en V1 en el patrón de Brugada.

\Needspace{8\baselineskip}\section*{08. Síndrome de Brugada y derivación V1}
\textbf{Tiempo orientativo: }30 segundos.

V1 es una derivación precordial derecha. En el esquema de la izquierda, la curva azul sirve de referencia y la roja ilustra el patrón de Brugada tipo 1. El QRS terminal es el último tramo del complejo QRS, justo antes del punto J, donde comienza el ST. En la curva roja, el punto J está elevado, el ST desciende desde ahí y la onda T se invierte. A la derecha se pueden comparar los tres patrones en las precordiales derechas. El simulador modifica por separado la parte final del QRS, el ST y la onda T.

\textbf{Transición.} A continuación veremos las dos formas de aprender a partir de las imágenes.

\Needspace{8\baselineskip}\section*{09. Redes neuronales convolucionales}
\textbf{Tiempo orientativo: }20 segundos.

Una convolución recorre la imagen con filtros cuyos coeficientes aprende la red. Un filtro puede resaltar, por ejemplo, los bordes del trazado. El pooling resume las activaciones y reduce su resolución espacial. La capa de clasificación usa esas características para producir una predicción. Durante el entrenamiento, el error respecto a las etiquetas guía el ajuste de los pesos.

\textbf{Transición.} En ICL, las etiquetas se incorporan como ejemplos en la consulta.

\Needspace{8\baselineskip}\section*{10. Aprendizaje en contexto}
\textbf{Tiempo orientativo: }25 segundos.

En el aprendizaje en contexto, el modelo recibe instrucciones, ejemplos etiquetados y una consulta nueva. El esquema ilustra el mecanismo con una tarea de texto. En este TFG se aplica a imágenes de ECG y los pesos se mantienen fijos durante la consulta.

\textbf{Transición.} Además de cómo aprende el modelo, importa si los ECG que recibe se parecen a los de la prueba.

\Needspace{8\baselineskip}\section*{11. Adaptación de dominio}
\textbf{Tiempo orientativo: }20 segundos.

El esquema de TarmacView utiliza la inspección de pavimentos para ilustrar el cambio entre los datos de entrenamiento y los de aplicación. En ECG, Gliner y colaboradores estudiaron un cambio parecido al pasar de trazados escaneados a fotos de móvil. En este TFG, el origen es el simulador y el destino son los registros del HUCA. La adaptación intenta reducir el efecto de ese cambio sobre el modelo.

\textbf{Transición.} Estos son los dos conjuntos con los que se estudia ese cambio.

\Needspace{8\baselineskip}\section*{12. Conjuntos de datos del simulador y del HUCA}
\textbf{Tiempo orientativo: }25 segundos.

Los dos conjuntos aportan imágenes de un latido de V1. El simulador genera \SimPatients{} casos y sus etiquetas describen rasgos morfológicos. El HUCA aporta \HucaPatients{} pacientes con una etiqueta clínica por persona. Los trazados que se muestran son ejemplos independientes.

\textbf{Transición.} Primero describo cómo se generan los casos simulados.

\Needspace{8\baselineskip}\section*{13. Simulador QRS/ST}
\textbf{Tiempo orientativo: }35 segundos.

El simulador construye cada latido sumando gaussianas. Su amplitud, posición y anchura permiten modificar regiones de las ondas. Sayadi y colaboradores utilizaron componentes gaussianos en un modelo dinámico. El simulador de este trabajo suma directamente los componentes. Las etiquetas se asignan a partir de los rasgos que aparecen en el latido generado.

\textbf{Transición.} Al variar sus componentes se construyen distintas familias de trazados.

\Needspace{8\baselineskip}\section*{14. Familias del simulador}
\textbf{Tiempo orientativo: }25 segundos.

Las familias incluyen un trazado de referencia y distintas combinaciones de cambios en el QRS, el ST y la onda T. Cada familia permite generar ejemplos con una morfología distinta.

\textbf{Transición.} Veamos cuántas veces aparece cada rasgo en el conjunto generado.

\Needspace{8\baselineskip}\section*{15. Distribución de etiquetas en el simulador}
\textbf{Tiempo orientativo: }25 segundos.

Los rasgos morfológicos pueden coexistir. La regla que exige los tres produce \SimPositive{} positivos y \SimNegative{} negativos entre los \SimPatients{} casos simulados. Este desequilibrio motiva el uso de métricas que describan el rendimiento en ambas clases y en la clase positiva.

\textbf{Transición.} El segundo conjunto aporta la evaluación sobre registros clínicos.

\Needspace{8\baselineskip}\section*{16. Conjunto clínico del HUCA}
\textbf{Tiempo orientativo: }25 segundos.

La campaña evaluó \HucaPatients{} pacientes, con \HucaPositive{} etiquetas positivas y \HucaNegative{} negativas. Según la memoria, se excluyeron casos que los colaboradores clínicos aconsejaron dejar fuera por dudas de clasificación. En el código, los 46 casos excluidos son los que tenían a la vez la etiqueta positiva de Brugada y el indicador de patrón basal patológico en la copia de trabajo. La copia pública de PhysioNet tiene etiquetas distintas. Los recuentos de la diapositiva corresponden a la muestra utilizada en el TFG.

\textbf{Transición.} Con ambos conjuntos definidos, el esquema sitúa los dos enfoques.

\Needspace{8\baselineskip}\section*{17. Diseño experimental}
\textbf{Tiempo orientativo: }30 segundos.

El esquema reúne dos formas de utilizar las mismas entradas. La CNN ajusta sus pesos con imágenes etiquetadas. Gemma utiliza demostraciones dentro de la consulta. Los dos enfoques se comparan en el simulador, en transferencia a registros clínicos y con ejemplos del HUCA.

\textbf{Transición.} La primera implementación es la CNN supervisada con ResNet18.

\Needspace{8\baselineskip}\section*{18. Implementación de la CNN con ResNet18}
\textbf{Tiempo orientativo: }35 segundos.

La CNN utiliza ResNet18 preentrenada y ajusta la red completa. La imagen se convierte a gris, se redimensiona y se replica su intensidad en tres canales para respetar la entrada de ResNet18. El simulador requiere tres salidas morfológicas. El entrenamiento con HUCA utiliza una salida para la etiqueta clínica.

\textbf{Transición.} Grad-CAM permite inspeccionar una salida de la red entrenada.

\Needspace{8\baselineskip}\section*{19. Inspección de activaciones con Grad-CAM}
\textbf{Tiempo orientativo: }30 segundos.

Grad-CAM permite inspeccionar qué regiones de la imagen se asocian con una salida de la CNN. Aquí se muestra la salida correspondiente al rasgo RBBB mediante la imagen, el mapa de activación y su superposición. El análisis es cualitativo. El mapa muestra una asociación con la salida del modelo.

\textbf{Transición.} La otra implementación utiliza demostraciones en una consulta a Gemma.

\Needspace{8\baselineskip}\section*{20. Implementación del ICL con Gemma 4 E4B}
\textbf{Tiempo orientativo: }25 segundos.

E4B es una variante compacta de la familia Gemma 4 de Google. La consulta contiene instrucciones, demostraciones etiquetadas y la imagen que se quiere clasificar. Se prueban selecciones estándar y balanceadas por clase. El modelo devuelve etiquetas estructuradas para calcular las métricas.

\textbf{Transición.} Una vez implementados los dos enfoques, los comparamos en tres escenarios.

\Needspace{8\baselineskip}\section*{21. Escenarios de evaluación}
\textbf{Tiempo orientativo: }15 segundos.

Primero se evalúan los dos enfoques dentro del simulador. Después se mantienen los ejemplos simulados y se pasa a evaluar en el HUCA, donde también se prueban estrategias de adaptación. Por último, tanto los ejemplos como la evaluación proceden del HUCA.

\textbf{Transición.} El protocolo concreta cómo se reserva un paciente para la prueba.

\Needspace{8\baselineskip}\section*{22. Protocolo de evaluación por paciente: LOOCV}
\textbf{Tiempo orientativo: }35 segundos.

K es el número de casos de entrenamiento o demostraciones. Se evalúan cinco valores con tres semillas. Cada pliegue reserva un paciente para la prueba. La CNN clínica dispone además de ocho pacientes de validación para elegir época y umbral. La exactitud balanceada promedia sensibilidad y especificidad. F1 combina precisión y sensibilidad de la clase positiva.

\textbf{Transición.} Empezamos con el escenario en el que todos los datos son simulados.

\Needspace{8\baselineskip}\section*{23. Simulación}
\textbf{Tiempo orientativo: }5 segundos.

En el primer escenario, los ejemplos y los datos de prueba son simulados.

\Needspace{8\baselineskip}\section*{24. Entrenamiento y evaluación con datos simulados}
\textbf{Tiempo orientativo: }55 segundos.

En este escenario, tanto los ejemplos como la prueba proceden del simulador. La CNN se entrena con esos casos y Gemma los utiliza como demostraciones. El eje horizontal indica cuántos ejemplos recibe cada enfoque. A la izquierda se muestra la exactitud balanceada y a la derecha, F1. Las dos curvas muestran una mejora al aumentar la cantidad de ejemplos. Con treinta y dos, la exactitud balanceada es \SimCnnBA{} para la CNN y \SimVlmBA{} para Gemma con selección balanceada. Las barras indican variación entre semillas. El resultado describe el reconocimiento de las reglas morfológicas del simulador.

\textbf{Transición.} Las matrices permiten distinguir los errores de cada clase.

\Needspace{8\baselineskip}\section*{25. Matrices de confusión en el simulador}
\textbf{Tiempo orientativo: }30 segundos.

Las matrices comparan CNN y Gemma con treinta y dos ejemplos. Cada una agrega tres evaluaciones de los mismos casos. Gemma con selección balanceada produce menos falsos positivos en esta configuración. La matriz permite distinguir ese efecto de los errores sobre la clase positiva.

\textbf{Transición.} Ahora cambiamos los datos de prueba.

\Needspace{8\baselineskip}\section*{26. Cambio de dominio}
\textbf{Tiempo orientativo: }5 segundos.

En el segundo escenario se mantienen los ejemplos simulados y se evalúa en el HUCA.

\Needspace{8\baselineskip}\section*{27. CNN entrenada en el simulador y evaluada en el HUCA}
\textbf{Tiempo orientativo: }35 segundos.

Aquí se entrena la CNN con datos del simulador y se cambia el conjunto de prueba. La línea continua representa la evaluación simulada y la discontinua, la clínica. Con treinta y dos ejemplos la exactitud balanceada clínica queda en \TransferCnnLastBA{}. El cambio de dominio afecta tanto al aspecto del trazado como a las etiquetas.

\textbf{Transición.} Las variantes de adaptación intentan reducir esa pérdida de rendimiento.

\Needspace{8\baselineskip}\section*{28. Métodos de adaptación de dominio en la CNN}
\textbf{Tiempo orientativo: }40 segundos.

Las variantes de adaptación incorporan imágenes clínicas sin etiqueta, incluidas las que luego se evalúan. Las etiquetas de entrenamiento siguen procediendo del simulador. CORAL y MMD mejoran algunos puntos frente a la CNN base con el mismo K. El mejor resultado global sigue correspondiendo a la CNN sin adaptación. El panel derecho muestra el resultado de Gemma con ejemplos simulados.

\textbf{Transición.} F1 permite comprobar qué sucede con la detección de positivos.

\Needspace{8\baselineskip}\section*{29. F1 de la clase positiva al cambiar de dominio}
\textbf{Tiempo orientativo: }30 segundos.

El panel de F1 permite examinar la detección de positivos al cambiar de dominio. El rendimiento sigue siendo bajo en la clasificación clínica, tanto con las variantes de adaptación como con las demostraciones simuladas de Gemma. Esto ayuda a interpretar una especificidad alta cuando se predice casi siempre la clase negativa.

\textbf{Transición.} En Gemma también podemos variar el origen de las demostraciones.

\Needspace{8\baselineskip}\section*{30. Cambio de dominio en ICL con selección balanceada}
\textbf{Tiempo orientativo: }35 segundos.

En la selección balanceada se compara qué ocurre al cambiar el origen de las demostraciones de Gemma, cuyos pesos permanecen fijos. Ante consultas clínicas, los ejemplos del HUCA ayudan a clasificar otros registros del mismo conjunto. Con demostraciones del simulador, el rendimiento clínico se mantiene próximo al azar.

\textbf{Transición.} La selección estándar permite comprobar si el patrón se mantiene con otro reparto de clases.

\Needspace{8\baselineskip}\section*{31. Cambio de dominio en ICL con selección estándar}
\textbf{Tiempo orientativo: }20 segundos.

La selección estándar reproduce el mismo patrón general. Las consultas clínicas se benefician de ejemplos del HUCA, mientras que las demostraciones simuladas aportan poco al clasificar los registros clínicos.

\textbf{Transición.} En el último escenario también cambia el origen de los ejemplos.

\Needspace{8\baselineskip}\section*{32. Evaluación clínica}
\textbf{Tiempo orientativo: }5 segundos.

En el tercer escenario, tanto los ejemplos como la prueba proceden del HUCA.

\Needspace{8\baselineskip}\section*{33. Entrenamiento y demostraciones con ejemplos del HUCA}
\textbf{Tiempo orientativo: }70 segundos.

En el tercer escenario, la CNN se entrena directamente con imágenes del HUCA y Gemma recibe demostraciones del mismo conjunto. El paciente de prueba queda reservado en ambos casos. A la izquierda vemos la exactitud balanceada. A la derecha se muestra F1 para la clase positiva. F1 baja cuando el modelo deja casos positivos sin detectar o cuando da demasiadas falsas alarmas. Sin demostraciones, Gemma no resulta capaz de identificar positivos en la evaluación realizada. Con dos o cuatro ejemplos, la mejor CNN conserva una ligera ventaja. Desde ocho ejemplos, Gemma con selección balanceada obtiene mejores resultados. Con treinta y dos alcanza una exactitud balanceada de \RealVlmBA{}, frente a \RealCnnLastBA{} de la CNN balanceada, y un F1 de \RealVlmFone{}. La comparación se refiere a esta muestra y a las configuraciones evaluadas.

\textbf{Transición.} Las matrices detallan los aciertos y errores por clase.

\Needspace{8\baselineskip}\section*{34. Matrices de confusión en la evaluación clínica}
\textbf{Tiempo orientativo: }35 segundos.

Las matrices muestran las mejores configuraciones balanceadas de cada método. La CNN utiliza dieciséis ejemplos y Gemma, treinta y dos. Gemma obtiene una sensibilidad de \RealVlmSens{} y una especificidad de \RealVlmSpec{}. Esto deja sin identificar alrededor de una cuarta parte de los positivos. Las matrices se normalizan por clase real y agregan tres evaluaciones de los mismos pacientes.

\textbf{Transición.} Antes de interpretar el alcance clínico, conviene recordar cómo se trataron los datos.

\Needspace{8\baselineskip}\section*{35. Aspectos éticos y protección de datos}
\textbf{Tiempo orientativo: }15 segundos.

El HUCA entregó los registros anonimizados. Los pesos abiertos de Gemma 4 permitieron ejecutar la inferencia en un clúster privado, sin enviar ECG a un servicio externo. La aplicación que se plantea es el apoyo al triaje: ordenar qué ECG revisar primero bajo responsabilidad del personal sanitario.

\textbf{Transición.} Estas condiciones también delimitan el alcance del estudio.

\Needspace{8\baselineskip}\section*{36. Discusión y limitaciones}
\textbf{Tiempo orientativo: }45 segundos.

Los ECG proceden de un solo centro, el HUCA. El simulador etiqueta rasgos morfológicos y los datos clínicos aportan una etiqueta por paciente. La entrada se limita a un latido de la derivación precordial V1. Elegimos un modelo de pesos abiertos para ejecutar la inferencia en nuestro entorno y evitar enviar los ECG a servicios externos. En ECGBench, algunos modelos cerrados superaron a varios modelos abiertos de propósito general en ciertas tareas. El rendimiento depende del modelo y de la tarea, y PULSE, un modelo abierto especializado en ECG, superó a los modelos cerrados evaluados en varias tareas. \href{https://aimedlab.github.io/PULSE/}{Liu y colaboradores (2026).}

\textbf{Transición.} Con ese alcance, se pueden responder los objetivos iniciales.

\Needspace{8\baselineskip}\section*{37. Conclusiones}
\textbf{Tiempo orientativo: }40 segundos.

Aumentar el número de ejemplos mejora el rendimiento de ambos métodos en la tarea simulada. En el HUCA, Gemma con selección balanceada supera a la CNN a partir de ocho ejemplos. Con 32 ejemplos, sus exactitudes balanceadas son \RealVlmBA{} y \RealCnnLastBA{}, respectivamente. Al entrenar la CNN con 32 ejemplos simulados y evaluarla en el HUCA, el resultado queda en \TransferCnnLastBA{}, cerca del azar.

\textbf{Transición.} Los resultados orientan la siguiente evaluación del enfoque.

\Needspace{8\baselineskip}\section*{38. Continuación del trabajo}
\textbf{Tiempo orientativo: }25 segundos.

La continuación empieza por revisar las etiquetas con especialistas y ampliar la recogida de casos etiquetados. Se estudiarían otros rasgos morfológicos del ECG y se probarían otros métodos de adaptación de dominio. Por último, se evaluaría si el sistema ayuda a priorizar ECG sospechosos en el HUCA y después en otros centros sanitarios.

\textbf{Transición.} Con esto termino la exposición.

\Needspace{8\baselineskip}\section*{39. Cierre y preguntas}
\textbf{Tiempo orientativo: }10 segundos.

Gracias por su atención. Quedo a disposición del tribunal para las preguntas.

\clearpage\section*{Referencias y preguntas previsibles}
\Needspace{7\baselineskip}\subsection*{Referencias clínicas y estado del arte}
Las referencias de contexto respaldan la motivación clínica y la procedencia de los datos y trazados. Los antecedentes corresponden al análisis de señales, el análisis de imágenes, el cambio de dominio y el aprendizaje con pocos ejemplos. Los créditos y las licencias de las imágenes publicadas se detallan en el proyecto LaTeX.

\Needspace{7\baselineskip}\subsection*{Referencias de métodos y métricas}
Estas fuentes identifican los métodos que aparecen en el desarrollo o los resultados, junto a la ficha del modelo y la definición de F1. Los valores experimentales proceden de los agregados canónicos del proyecto.

\section*{Preguntas que conviene preparar}
\begin{description}[style=nextline]
\item[¿La comparación utiliza el mismo presupuesto de etiquetas?] El valor de K es común, pero la CNN con HUCA dispone además de ocho pacientes de validación. Se debe declarar esta diferencia al explicar el protocolo.
\item[¿Por qué el simulador no transfiere bien?] Cambian tanto la apariencia de las imágenes como el significado de la etiqueta. El experimento no separa ambos efectos.
\item[¿La ventaja de Gemma demuestra significación estadística?] No. Las barras representan la desviación entre tres semillas y no sustituyen un contraste estadístico ni una validación externa.
\item[¿Podría usarse ya para descartar pacientes?] La mejor configuración todavía omite positivos. La utilidad como apoyo a la priorización necesita evaluación clínica, externa y prospectiva.
\item[¿Qué demuestra Grad-CAM?] Muestra regiones asociadas a una salida. No demuestra por sí solo un mecanismo causal ni que la CNN reconozca un criterio clínico.
\item[¿Qué ocurre sin demostraciones?] Con estas instrucciones, Gemma no identifica positivos. Es un resultado de la configuración evaluada y no una afirmación universal sobre todos los modelos o prompts.
\end{description}
\clearpage\section*{Fuentes para la preparación}
\begin{itemize}
\item Manual de Identidad Corporativa de UniOvi, apartado 6.1, páginas 57--58. \url{https://www.uniovi.es/conocenos/uniovi/identidad}
\item Página de TFG de la EPI. \url{https://epigijon.uniovi.es/estudiantes/tfg}
\item Copia pública de las directrices de entrega y defensa de la EPI, utilizada únicamente como orientación temporal. \url{https://expydoc.com/doc/6397600/directrices-para-los-estudiantes}
\item MIT Communication Lab, Slide Design. \url{https://mitcommlab.mit.edu/aeroastro/commkit/slide-design/}
\item Beamer User Guide, capítulo 5. \url{https://ctan.org/pkg/beamer}
\end{itemize}
\end{document}
